\documentclass[conference]{IEEEtran}
\IEEEoverridecommandlockouts
\usepackage{cite}
\usepackage{amsmath,amssymb,amsfonts}
\usepackage{graphicx}
\usepackage{textcomp}
\usepackage{xcolor}
\usepackage{booktabs}
\usepackage{algorithm}
\usepackage{algpseudocode}

\newcommand{\tbf}{\textcolor{red}{[to be filled in]}}
\newcommand{\tbc}[1]{\textcolor{blue}{#1}} % values to confirm against the registered design / network files
\newcommand{\phbox}[1]{\fbox{\parbox[c][#1][c]{0.95\columnwidth}{\centering\Large \textcolor{red}{[to be filled in]}}}}
% Implementation source for Sec. IV: E:/agent-work/c4-coopnative docs/paper_results/COOPMWM_DESIGN.md (items cited as D#), EXACT_GRADE_README.md, tools/coopmwm/{exact_grade,exploration,r_matching}.py

\begin{document}

\title{World-Model-Guided Exploration for Learning\\Multi-CAV Cooperation in Mixed Traffic}

\author{\IEEEauthorblockN{Author}}

\maketitle

\begin{abstract}
Connected and automated vehicles (CAVs) can reduce merge congestion by coordinating, for example by slowing together to meter inflow or by opening gaps for merging vehicles. In our experiments, a MAPPO baseline did not exhibit statistically significant positive synergy between CAVs, although exact simulator replays show that coordinated maneuvers can reduce delay relative to MAPPO. % COOPAUDIT.md; COOPPLAN.md
We propose World-Model-Guided Exploration (WGE) to improve the discovery of coordinated behavior. During training, a physics-informed traffic model searches for joint acceleration plans for groups of CAVs around the current policy; candidate plans are checked by short exact simulator branches, and the executed plans and their measured outcomes are used for reinforcement learning and selective imitation. The world model guides data collection but supplies no rewards. We evaluate the final policy with guidance disabled on 96 held-out scenarios: \tbf.
\end{abstract}

\begin{IEEEkeywords}
Connected automated vehicles, mixed traffic, multi-agent reinforcement learning, world models, cooperative control.
\end{IEEEkeywords}

%% =====================================================================
\section{Introduction}

At highway on-ramps, CAVs can improve traffic not only by driving well individually but by coordinating: several CAVs can slow down together to meter inflow, open gaps for merging vehicles, or hold speed across lanes so that the merge does not break down. Methods that select from a fixed library of such maneuvers restrict the joint actions that can be considered. We study whether coordinated behavior can be learned directly in the space of joint acceleration profiles.

Multi-agent reinforcement learning (MARL), in particular MAPPO~\cite{yu_mappo2022}, is widely used for CAV control~\cite{chen_merge2021,wu_flow2022}. Coordinated maneuvers are hard for it to discover: they help only in specific traffic states, their effect appears minutes later in a team reward shared by all vehicles, and independent per-vehicle exploration noise rarely produces a sustained joint maneuver. Consistent with this, the MAPPO baseline we use did not exhibit statistically significant positive synergy in the evaluated scenarios, while a planner with exact simulator look-ahead reduced delay relative to it (Sec.~VI-A).

We propose World-Model-Guided Exploration (WGE) to improve the discovery of coordinated CAV behavior. During training, a physics-informed traffic model searches for joint acceleration plans around the current policy. A short list of candidates is checked by exact simulator branches, the selected plan is executed in SUMO, and its measured outcome supports policy learning through reinforcement learning and selective imitation. The world model therefore guides data collection without supplying synthetic rewards. We evaluate whether the resulting policy reduces delay and exhibits positive counterfactual synergy when guidance is disabled.

Our contributions are:
\begin{itemize}
\item \textbf{Counterfactual evaluation of cooperation.} A together-versus-alone measure of CAV cooperation computed from exact simulator replays, applied to a planner with exact look-ahead, to MAPPO, and to learned policies.
\item \textbf{WGE.} A training method in which a physics-informed world model proposes coordinated group plans around the current policy, exact simulator branches select and grade them, and the policy learns from executed outcomes with PPO and advantage-filtered imitation.
\item \textbf{Empirical evidence.} On 96 held-out scenarios with guidance disabled, comparisons with MAPPO, an equal-budget unguided control and a random-guide ablation, together with an analysis of learned group behavior: \tbf.
\end{itemize}

%% =====================================================================
\section{Related Work}

\textbf{Learning CAV cooperation.} MARL has been applied to on-ramp merging and mixed-traffic control~\cite{chen_merge2021,guo_mappopis2024,wu_flow2022}. These methods differ in what they constrain: neighbor-averaged rewards that limit credit to local outcomes~\cite{chen_merge2021}; shared intentions communicated between agents~\cite{guo_mappopis2024}; altruistic or social-value objectives that change what agents optimize~\cite{toghi_social2022,schwarting_svo2019}; control restricted to selected road segments~\cite{vinitsky_lagrangian2018}; and hand-designed controllers~\cite{lee_circles2024}. Cooperation in these works is usually evaluated through aggregate traffic metrics and examples. We add an exact together-versus-alone measure and a training method that explores joint plans directly.

\textbf{Credit and guided exploration.} Difference rewards~\cite{wolpert_tumer2002,castellini_drreinforce} and COMA~\cite{foerster_coma2018} estimate an agent's contribution to a shared reward; our synergy measure applies the same counterfactual logic to groups, using exact replays. Cooperative exploration methods~\cite{liu_cmae2021} encourage joint exploration without a predictive model. Guide-based methods let a prior policy steer early training~\cite{uchendu_jsrl2023,torrey_advising2013}; in WGE the guide is a model-based planner whose use adapts to measured outcomes.

\textbf{Planning and world models.} Expert iteration and AlphaZero/MuZero improve a policy with look-ahead search~\cite{anthony_exit2017,silver_alphazero2018,schrittwieser_muzero2020}; MPPI and related planners combine sampling-based search with learned models~\cite{williams_mppi2017,hansen_tdmpc2_2024,levine_gps2013}. Multi-agent world models~\cite{egorov_mamba2022,zhang_marie2024} and a physics-informed model for offline mixed-traffic control~\cite{liu_physics_wm2026} train policies on imagined rollouts. WGE uses the model only to propose plans; all learning signals come from executed simulator trajectories.

%% =====================================================================
\section{Problem Formulation}

\subsection{Setting}
We study a highway segment with an on-ramp in SUMO~\cite{lopez_sumo2018}, with human-driven vehicles (HDVs) and CAVs. HDVs follow IDM~\cite{treiber_idm2000} car following and LC2013 lane changing with individual, unobserved parameters. Every 1\,s each CAV outputs a longitudinal acceleration that passes through a safety filter; lane changes are left to SUMO; the simulator step is 0.5\,s. Each CAV observes its own state and nearby vehicles; a roadside unit (RSU) observes the merge and ramp zones and entrance counts and broadcasts a learned 4-dimensional message. % D9
No agent observes driver parameters or future arrivals.

\subsection{Objective and reward}
Team delay over $[t_0,t_1]$ is
\begin{equation}
J[t_0,t_1] = \sum_{t=t_0}^{t_1}\Big[\sum_{v}\frac{v^{\mathrm{free}}_v - v_{v,t}}{v^{\mathrm{free}}_v} + N^{\mathrm{wait}}_t\Big]\Delta t ,
\end{equation}
% exact normalisation to be matched to the comparator reducer in tools/coopnative/comparator.py
the speed shortfall of vehicles on the road plus vehicles waiting to enter (veh-s). Lower is better. The shared per-step reward is $r_t = -(J[0,t+1]-J[0,t])$. % D4
Differences between policies are reported as policy minus reference; negative values mean less delay.

\subsection{Counterfactual cooperation}
For a group $G$ at state $s$, let $J_\varnothing$, $J_i$ and $J_G$ be the team delay over a common window when no member deviates from a reference policy, when only member $i$ executes its part of a joint plan, and when all members execute it. The member and group effects are $D_i = J_\varnothing - J_i$ and $D_G = J_\varnothing - J_G$, and the synergy is
\begin{equation}
C_G = D_G - \sum_{i\in G} D_i .
\end{equation}
Positive $C_G$ alone does not imply that the group helps: individual effects can be negative. We call a group behavior \emph{beneficial cooperation} if $D_G>0$ and $C_G>0$, and report both quantities. Computing them requires $|G|+2$ exact runs from an identical prefix; SUMO is deterministic given its state. Each audit states its reference policy (Sec.~VI). These replays are used for evaluation only.

\subsection{Goal}
We seek a policy that reduces $J$ relative to MAPPO when deployed without guidance, and whose group behaviors show beneficial cooperation.

%% =====================================================================
\section{Method}

\begin{figure}[t]
\centering
\phbox{4.2cm}
\caption{WGE overview.}
\label{fig:method}
\end{figure}

\subsection{Policy and critic}
Each CAV's action is $a_{i,t}=a^{\mathrm{M}}_{i,t}+u_{i,t}$, where $a^{\mathrm{M}}$ is the command of a frozen MAPPO actor~\cite{yu_mappo2022} and $u_{i,t}$ is a learned residual; $u=0$ reproduces MAPPO exactly. The residual network uses local telemetry, sensed neighbors and the RSU message. Exploration noise is an AR(1) process ($\phi=e^{-1/5}$ per step, $\sigma=0.5$\,m/s$^2$) with a component shared by nearby CAVs (correlation $0.5$); the PPO likelihood uses the conditional Gaussian density of each action given the previous noise state. % D11
A centralized critic pools CAV features with 12 RSU zone observations. We use PPO with 4 epochs, Adam ($3\times10^{-4}$), clip 0.2, a KL limit of 0.045, $\gamma=e^{-1/600}$ and $\lambda=e^{-1/1200}$ per 1\,s step. % D12, D13

\subsection{World model}
The world model rolls traffic forward from what the CAVs and RSU sense (position, speed, lane and type of each sensed vehicle, and entrance queues) in 0.5\,s steps. CAVs follow the evaluated plan through the same safety filter; HDVs follow IDM,
\begin{equation}
\dot v = a_{\max}\Big[1-\big(\tfrac{v}{v_0}\big)^{\delta}-\big(\tfrac{s^*}{s}\big)^2\Big],\quad
s^* = s_0 + vT + \tfrac{v\,\Delta v}{2\sqrt{a_{\max} b}},
\end{equation}
and SUMO's lane-change and merge rules. Each HDV's parameters $(v_0,T,s_0,a_{\max},b)$ are drawn from a posterior inferred from that driver's last 60\,s of sensed motion, shrunk toward a learned population prior when history is short; the encoder is trained on prediction error only, without driver labels. Arrival rates are inferred from entrance counts. A forecast averages $K$ rollouts with independent draws of these unknowns; competing plans share the same $K$ draws. % D17, D18

\subsection{Group plans and search}
At a decision step, CAVs anywhere on the road (all lanes and the ramp) are chained into groups of 2--6 members if they are within 100\,m of one another. % D19
A plan assigns each member an acceleration profile of duration $D\in\{30,60,\dots,180,240,300\}$\,s, parameterized by at most 10 piecewise-linear knots and bounded by the vehicle limits. % D69
The search is centered on the current policy. The model first rolls out the policy's own trajectory for the group; a cross-entropy/MPPI search~\cite{williams_mppi2017} then samples $N=\max(32, 20|G|)$ joint perturbations of that trajectory per iteration, smooth in time and correlated across members, screens them with $K=4$ rollouts over 300\,s, keeps the best $\lceil N/4\rceil$ and resamples, for 3 iterations. Eight finalists are re-scored with $K=8$ over $H=\min(600,\max(300,D+120))$\,s. The search uses no hand-designed maneuvers. The top four finalists with positive predicted saving over MAPPO form the shortlist. % D21, D22
The perturbation covariance is updated from plans that later receive positive simulator grades (Sec.~IV-E); maneuver families are used only to label results afterward.

\subsection{Selection, execution and grading}
At an eligible decision step, a guided group follows a plan with probability $\beta$ (the trust, Sec.~IV-F), drawn before any further computation. If it follows, the shortlist and the policy's own continuation are evaluated by five exact SUMO branches from the identical prefix with common random numbers, over a common window ending $\max D+120$\,s later, and the best branch is executed; if that branch is the policy itself, no plan is executed. % D42, D64
Executed plans run to completion or to the episode end, and are recorded as truncated in the latter case.

Each executed plan $p$ starting at $t_0$ with duration $D$ receives an exact grade
\begin{equation}
g_p = J^{\text{without}}[t_0,t_e] - J^{\text{with}}[t_0,t_e],
\end{equation}
with $t_e=\min(t_0+D+120,\,1200)$\,s,
where ``with'' is the executed episode and ``without'' is an exact rerun from the identical prefix in which the plan's members follow the policy's own stochastic actions (same noise draws) and every other executed plan in the window is replayed as recorded. Plans that share members with another plan in the window are excluded from grading. % EXACT_GRADE_README.md
These training branches (selection and grading) are counted in the simulation budget (Sec.~V-C).

\subsection{Learning}
The actor loss is the standard clipped PPO objective over steps at which the policy chose its own action, including steps where a plan was considered but not executed. Steps inside an executed plan contribute no PPO ratio. The critic is fit to the observed returns on all steps. The advantages of policy steps are computed by GAE on the observed trajectory, so for steps preceding an intervention they include the effect of the executed plan; we report this as a limitation of the current estimator. % D27, D28
In addition, executed plans with positive grade contribute a weighted imitation loss: the negative log-likelihood of each member's plan accelerations under the policy, with weight 1 until 20 grades are available and $w_p=\min(20, \exp(g_p/\tau))$ thereafter, where $\tau=\max(1, \mathrm{RMS}(g_{1..20}))$ is fixed from the first 20 grades. Only plans with $g_p>0$ contribute; the full positive buffer is used at every update. This is a filtered variant of advantage-weighted regression~\cite{peng_awr2019}. % D28

\subsection{Trust}
For each bucket of traffic regime (merge-dense or ordinary, defined from observable CAV-pair geometry and CAV/density bands; \tbc{exact definition to add}), let $m$ and $[L,U]$ be the mean and 95\% world-clustered bootstrap interval of the exact grades so far, and $e=\max(|m-L|,|U-m|)$. After each update,
\begin{equation}
\beta \leftarrow \min\Big(\beta,\ \frac{\max(0,U)}{|m|+e}\Big),
\end{equation}
with empirical-Bayes shrinkage of bucket estimates toward the pooled estimate; $\beta$ starts at 1, requires grades from at least two worlds, and does not increase. Training pauses if at least 30 grades have $U<0$. % D65, D66; exact_grade.trust_evidence
This rule reduces guidance under uncertain or negative evidence but does not force $\beta$ to zero. At evaluation, guidance is disabled.

\subsection{Random-guide ablation}
The random guide (R) uses the same policy-centered perturbation distribution, exact selection branches, grading, imitation and trust rule, but its candidates are sampled without world-model scoring. Its executed plans are matched to those of WGE update by update: by regime first, then by group size within one member, with the same durations; shortfalls are reported. % D16; r_matching.py
WGE versus R therefore isolates the contribution of model-based candidate generation; WGE versus the unguided control (A) measures the full method, whose learning objective additionally includes the imitation term.

\begin{algorithm}[t]
\caption{WGE training (one update)}
\label{alg:wge}
\small
\begin{algorithmic}[1]
\For{each of 24 training episodes}
  \For{each 1\,s decision step}
    \If{an eligible group exists and a plan is requested (prob.\ $\beta$)}
      \State shortlist $\gets$ policy-centered world-model search
      \State select best of shortlist $\cup$ \{policy\} by 5 exact branches
      \State execute selected plan to completion (or episode end)
    \Else \State act with the policy (residual on MAPPO)
    \EndIf
  \EndFor
\EndFor
\State grade executed plans exactly (Eq.~5)
\State update actor (PPO on own steps), critic (all steps), imitation (plans with $g_p>0$)
\State update $\beta$ (Eq.~6), search covariance and plan library
\end{algorithmic}
\end{algorithm}

%% =====================================================================
\section{Experimental Setup}

\subsection{Road and traffic}
SUMO~1.27.1 simulates a highway segment with an on-ramp (\tbc{length and lane counts to confirm}). Scenarios combine mainline demand $\{1900,2400,2700\}$\,veh/h with CAV penetration $\{5,10,20,40\}\%$ (12 traffic types). Episodes last 1200\,s: 300\,s warm-up, then 900\,s of control. CAV sensing range: \tbc{100\,m}.

\subsection{Compared policies}
\textbf{Native}: SUMO default driving. \textbf{MAPPO}: speed-control MAPPO trained for 169 updates ($\approx$4.7\,M controlled simulation-seconds). % COOPNATIVE_DESIGN.md: 169x46x600
\textbf{Control (A)}: MAPPO plus the same residual network trained with PPO, without guidance. \textbf{WGE (B)}: Sec.~IV. \textbf{Random guide (R)}: Sec.~IV-G. A, B and R start from the same MAPPO policy and use the same network, training scenarios and episodes per update.

\subsection{Training and budget}
Learned policies train on 72 scenarios (6 per traffic type) separate from all test scenarios. Each update uses 24 episodes, half from near-capacity, high-penetration types (2400--2700\,veh/h, 20--40\% CAVs). A, B and R are trained for 24 updates with one training seed each; A is additionally continued to 48 updates as a secondary result. Each B and R update requires at least 10 executed plans. Simulation budgets, including selection and grading branches, and world-model GPU time are reported for every arm: \tbf. % D44, D58
Progress is tracked on 12 separate scenarios that are never used for model selection.

\subsection{Evaluation}
All final results use 96 held-out scenarios (8 per traffic type) registered before training and evaluated once with the final (update-24) checkpoint of each policy, with guidance disabled. % D59, D39
\begin{itemize}
\item \textbf{Paired branch test.} At $t=300,500,700,900$\,s every policy continues for 300\,s from the same state, reached by MAPPO driving up to that time. The outcome is paired 300\,s team delay. % D61
\item \textbf{Near-capacity subset.} The same comparison on the 32 near-capacity, high-penetration scenarios, selected in advance from earlier evidence.
\end{itemize}
Intervals are 95\% bootstrap intervals with 10{,}000 resamples of whole scenarios, keeping all four branches of a scenario together. Holm correction is applied in two families, one per comparison (B$-$MAPPO and B$-$A), each over the two tests above at $\alpha=0.05$; B$-$R and A$-$MAPPO are secondary. % D67
Secondary outcomes are full-episode delay over 600 and 900\,s, per-type results, percentage reduction and safety (collisions, time-to-collision events, hard braking). Beneficial cooperation (Sec.~III-C) is evaluated for a random sample of executed plans and for group behaviors detected in the final policy.

%% =====================================================================
\section{Results}

\subsection{Cooperation above MAPPO (E1)}
\emph{Exact look-ahead planner.} Every 60\,s, a planner replays candidate group maneuvers exactly from the current state and executes one only if it beats MAPPO (reference: MAPPO). On 36 new scenarios covering all traffic types, it reduces delay relative to MAPPO by 648\,veh-s per scenario over 600\,s (95\% CI [210, 1174]) and by 1584\,veh-s over 900\,s ([588, 2765]). % COOPPLAN.md Pooled36: J600 +648.407 [209.699,1173.616]; J900 +1583.821 [588.246,2765.414]
Of 52 executed group maneuvers, 24 have $C_G>0$, including one whose joint execution reduced delay although each member's part alone increased it; one cross-lane wall had $C_G=+182$\,veh-s. % COOPPLAN.md H300 amendment: C_G +181.571
Mean synergy across maneuvers is positive but not significant.

\emph{MAPPO.} In exact replays of 21 scenarios with pairs of nearby CAVs, with native driving as the reference for each member, MAPPO's synergy is $+68$\,veh-s over 600\,s (95\% CI [$-1179$, $+1154$]); we did not detect significant positive synergy. In point estimate, 825 of its 893\,veh-s group effect comes from individual effects. % COOPAUDIT.md Team J600 mappo

\emph{Rarity.} Among 300 training states, our maneuver search found a helpful option in 15\% (95\% CI [10\%, 20\%]); averaged over all states, each maneuver type increased delay by 1000--1700\,veh-s, while the best cases saved up to 17{,}000\,veh-s. % COOPBANDIT.md: 45/300 [10%,20.33%]; mean effects -1176/-1153/-1736/-1022; best +14992/+9917/+6338/+17154
A learned maneuver selector never acted, and a policy distilled from the planner agreed with it on 85\% of decisions but on none of the 45 helpful ones. % COOPBANDIT.md; DISTILL.md (teacher agreement 84.67%, helpful-teacher agreement 0/45)

\begin{figure}[t]
\centering
\phbox{4.2cm}
\caption{Cooperation above MAPPO: (a) planner saving per scenario; (b) $C_G$ and $D_G$ of planner maneuvers; (c) MAPPO's group effect decomposition.}
\label{fig:opportunity}
\end{figure}

\subsection{World-model selection quality (E2)}
On 12 states with exact labels (58 candidate pairs per horizon), the model orders pairs of maneuvers correctly in 65.5\% of cases at 120\,s (Wilson 95\% interval [53\%, 76\%]) and 75.9\% at 300\,s ([64\%, 85\%]); pairs within a state are not independent. % BDRANK.md: 58 pairs/horizon; intervals computed from n=58, ignore clustering
Ranking by predicted delay outperforms ranking by predicted breakdown probability, which rarely separates plans (ties in 112 of 120 cases). % BDRANK.md
More directly, we report the exact grade of executed plans (Eq.~5) relative to the current policy, by update and by predicted rank: \tbf.

\subsection{Main comparison (E3)}
\tbf

\begin{table}[t]
\caption{Delay difference relative to each reference (veh-s, 95\% CI; negative is better).}
\label{tab:main}
\centering
\phbox{3.2cm}
\end{table}

\begin{figure}[t]
\centering
\phbox{4.2cm}
\caption{Training progress, trust $\beta$ and executed plans per update.}
\label{fig:curves}
\end{figure}

\subsection{Learned group behavior (E4)}
\tbf

\begin{figure}[t]
\centering
\phbox{4.2cm}
\caption{A group behavior of the final WGE policy and its counterfactual replays.}
\label{fig:behavior}
\end{figure}

%% =====================================================================
\section{Discussion and Limitations}

In our experiments, exact look-ahead found coordinated maneuvers that reduce delay relative to MAPPO, while the MAPPO baseline showed no significant synergy. \tbf.

\textbf{Limitations.} Results are from simulation with SUMO's driver models on one on-ramp; other layouts and real drivers may differ. Each policy has one training run, so variation across training seeds is not measured. MAPPO is our implementation. WGE's training uses exact simulator branches for selection and grading, which require a resettable simulator; the deployed policy does not. The PPO advantages of steps preceding an intervention include the intervention's effect. WGE and the control differ in their learning objective as well as in data collection; the random guide shares WGE's objective. The trust rule does not increase $\beta$ once reduced. The world model is fitted to this road and is less accurate for long plans.

\textbf{Future work} includes multiple training seeds, an unguided control with comparable self-imitation, intervention-aware advantage estimation, other road layouts and real trajectories.

%% =====================================================================
\section{Conclusion}

We studied whether coordinated CAV behavior can be learned in the space of joint acceleration plans. Exact replays showed coordinated maneuvers that reduce delay relative to a MAPPO baseline without detectable synergy. WGE uses a physics-informed world model to propose group plans around the current policy, checks and grades them in the simulator, and learns from the executed outcomes. \tbf.

\section*{Acknowledgment}
To be added.

\bibliographystyle{IEEEtran}
\bibliography{refs}

\end{document}
